\subsection{GRADED MODULES OF GRADED HOMOMORPHISMS}

We suppose in this no. that the monoid \(\Delta\) is a group. Let A be a graded ring of type \(\Delta\) and M, N two graded left (for example) A-modules of type \(\Delta\) . Let H \(_{\lambda}\) denote the additive group of graded homomorphisms of degree \(\lambda\) of M into N (no. 2); in the additive group Hom \(_{A}\) (M, N) of all homomorphisms of M into N (with the non-graded A-module structures) the sum (for \(\lambda \in \Delta\) ) of the H \(_{\lambda}\) is direct. For, if there is a relation \(\sum_{\lambda} u_{\lambda} = 0\) with \(u_{\lambda} \in \mathrm{H}_{\lambda}\) for all \(\lambda\) , it follows that \(\sum_{\lambda} u_{\lambda}(x_{\mu}) = 0\) for all \(\mu\) and all \(x_{\mu} \in \mathrm{M}_{\mu}\) . As the elements of \(\Delta\) are cancellable, \(u_{\lambda}(x_{\mu})\) is the homogeneous component of \(\sum_{\lambda} u_{\lambda}(x_{\mu})\) of degree \(\lambda + \mu\) ; hence \(u_{\lambda}(x_{\mu}) = 0\) for every ordered pair ( \(\mu, \lambda\) ) and every \(x_{\mu} \in \mathrm{M}_{\mu}\) , which implies \(u_{\lambda} = 0\) for all \(\lambda \in \Delta\) . We shall denote (in this paragraph) by Homgr \(_{A}\) (M, N) the additive subgroup of Hom \(_{A}\) (M, N) the sum of the H \(_{\lambda}\) and we shall call it the additive group of graded A-module homomorphisms of M into N. Let C be the centre of A, which is a graded subring (no. 3, Corollary to Proposition 5); for the canonical C-module structure on \(\mathrm{Hom}_{\mathbf{A}}(\mathbf{M},\mathbf{N})\) (§ 1, no. 14, Remark 1), \(\mathrm{Homgr}_{\mathbf{A}}(\mathbf{M},\mathbf{N})\) is a submodule and the graduation \((\mathbf{H}_{\lambda})\) is compatible with the C-module structure: for, if \(c_{\nu}\in \mathbf{C}\cap \mathbf{A}_{\nu}\) , \(x_{\mu}\in \mathbf{N}_{\mu}\) and \(u_{\lambda}\in \mathbf{H}_{\lambda}\) , then by definition \((c_{\nu}u_{\lambda})(x_{\mu}) = c_{\nu}.u_{\lambda}(x_{\mu})\subset \mathbf{N}_{\lambda +\mu +\nu}\) and hence \(c_{\nu}u_{\lambda}\in \mathbf{H}_{\lambda +\nu}\) .

Let \(\mathbf{M}'\) and \(\mathbf{N}'\) be two graded left A-modules of type \(\Delta\) and \(u':\mathbf{M}'\to\mathbf{M}\) , \(v':\mathbf{N}\to\mathbf{N}'\) graded homomorphisms of respective degrees \(\alpha\) and \(\beta\) . Then it is immediate that \(\operatorname{Hom}(u',v')\colon w\mapsto v'\circ w\circ u'\) maps \(\operatorname{Homgr}_{\mathbf{A}}(\mathbf{M},\mathbf{N})\) into \(\operatorname{Homgr}_{\mathbf{A}}(\mathbf{M}',\mathbf{N}')\) and that its restriction to \(\operatorname{Homgr}_{\mathbf{A}}(\mathbf{M},\mathbf{N})\) is a graded homomorphism into \(\operatorname{Homgr}_{\mathbf{A}}(\mathbf{M}',\mathbf{N}')\) of degree \(\alpha+\beta\) .

In particular \(\operatorname{Homgr}_{\mathbf{A}}(\mathbf{M}, \mathbf{M})\) is a graded subring of \(\operatorname{End}_{\mathbf{A}}(\mathbf{M})\) , which is denoted by \(\operatorname{Endgr}_{\mathbf{A}}(\mathbf{M})\) .

Remark. If M and N are graded left A-modules, \(\operatorname{Homgr}_{\mathbf{A}}(\mathbf{M}, \mathbf{N})\) is in general distinct from \(\operatorname{Hom}_{\mathbf{A}}(\mathbf{M}, \mathbf{N})\) . However these two sets are equal when M is a finitely generated A-module. For M is then generated by a finite number of homogeneous elements \(x_{i}\) ( \(1 \leqslant i \leqslant n\) ); let \(d(i)\) be the degree of \(x_{i}\) ; let \(u \in \operatorname{Hom}_{\mathbf{A}}(\mathbf{M}, \mathbf{N})\) and for all \(\lambda \in \Delta\) let \(z_{i,\lambda}\) denote the homogeneous component of \(u(x_{i})\) of degree \(\lambda + d(i)\) . We show that there exists a homomorphism \(u_{\lambda}: \mathbf{M} \to \mathbf{N}\) such that \(u_{\lambda}(x_{i}) = z_{i,\lambda}\) for all \(i\) . It suffices to prove that if \(\sum_{i} a_{i} x_{i} = 0\) with \(a_{i} \in \mathbf{A}\) for \(1 \leqslant i \leqslant n\) , then \(\sum_{i} a_{i} z_{i,\lambda} = 0\) for all \(\lambda \in \Delta\) (§ 1, no. 7, Remark). It can be assumed that each \(a_{i}\) is homogeneous of degree \(d'(i)\) such that \(d(i) + d'(i) = \mu\) for all \(i\) (no. 3, Remark 1); then \(\sum_{i} a_{i} u(x_{i}) = 0\) ; taking the homogeneous component of degree \(\lambda + \mu\) on the left-hand side, we obtain \(\sum_{i} a_{i} z_{i,\lambda} = 0\) , whence the existence of the homomorphism \(u_{\lambda}\) ; clearly moreover \(u_{\lambda}\) is graded of degree \(\lambda\) . Finally, \(u_{\lambda} = 0\) except for a finite number of values of \(\lambda\) , and \(u = \sum_{\lambda} u_{\lambda}\) by definition, which proves our assertion.

In particular, \(\operatorname{Homgr}_{\mathbf{A}}(\mathbf{A}_s, \mathbf{M}) = \operatorname{Hom}_{\mathbf{A}}(\mathbf{A}_s, \mathbf{M})\) for every graded left A-module M; moreover \(\operatorname{Hom}_{\mathbf{A}}(\mathbf{A}_s, \mathbf{M})\) has a graded left A-module structure (and not just a graded C-module structure), and it is immediate that with this structure the canonical mapping of M into \(\operatorname{Hom}_{\mathbf{A}}(\mathbf{A}_s, \mathbf{M})\) ( \(\S 1\) , no. 14, Remark 2) is a graded A-module isomorphism.

Similarly, \(\operatorname{Homgr}_{\mathbf{A}}(\mathbf{M}, \mathbf{A}_s)\) has a graded right A-module structure (and not only a graded C-module structure); it is called the graded dual of the graded A-module M and is denoted by \(\mathbf{M}^{*\mathbf{gr}}\) , or simply \(\mathbf{M}^*\) when no confusion results. If \(u: \mathbf{M} \to \mathbf{N}\) is a graded homomorphism of degree \(\delta\) , it follows from the above that the restriction to \(\mathbf{N}^{*\mathbf{gr}}\) of \(^t u = \operatorname{Hom}(u, 1_{\mathbf{A}_s})\) is a graded homomorphism of the graded dual \(\mathbf{N}^{*\mathbf{gr}}\) into the graded dual \(\mathbf{M}^{*\mathbf{gr}}\) , of degree \(\delta\) , called the graded transpose of \(u\) .

We sometimes consider on the graded dual \(M^{*gr}\) the graduation derived from the above with the aid of the isomorphism \(\lambda\mapsto-\lambda\) of \(\Delta\) (no. 1, Example 2) so that the homogeneous elements of degree \(\lambda\) in \(M^{*gr}\) are the graded linear forms of degree \(-\lambda\) on M (when A has the trivial graduation, these are the zero linear forms of the \(M_{\mu}\) of index \(\mu\neq\lambda\) ). Then, if \(u:M\to N\) is a graded homomorphism of degree \(\delta\) , \(u'\) becomes a graded homomorphism of degree \(-\delta\) .

Suppose A is commutative and graded of type \(\Delta\) and let M, N, P, Q be four graded A-modules of type \(\Delta\) . Then there are canonical graded homomorphisms of degree 0

\begin{enumerate}
\def\labelenumi{(\arabic{enumi})}
\setcounter{enumi}{2}
\tightlist
\item
\end{enumerate}

\[
\operatorname{Homgr} _ {\mathbf {A}} (\mathbf {M}, \operatorname{Homgr} _ {\mathbf {A}} (\mathbf {N}, \mathbf {P})) \rightarrow \operatorname{Homgr} _ {\mathbf {A}} (\mathbf {M} \otimes_ {\mathbf {A}} \mathbf {N}, \mathbf {P})\tag{4}
\]

\[
\operatorname{Homgr} _ {\mathbf {A}} (\mathbf {M}, \mathbf {N}) \otimes_ {\mathbf {A}} \mathrm{P} \rightarrow \operatorname{Homgr} _ {\mathbf {A}} (\mathbf {M}, \mathbf {N} \otimes_ {\mathbf {A}} \mathrm{P})
\]

\[
\operatorname{Homgr} _ {\mathbf {A}} (\mathbf {M}, \mathbf {P}) \otimes_ {\mathbf {A}} \operatorname{Homgr} _ {\mathbf {A}} (\mathbf {N}, \mathbf {Q}) \rightarrow \operatorname{Homgr} _ {\mathbf {A}} (\mathbf {M} \otimes_ {\mathbf {A}} \mathbf {N}, \mathbf {P} \otimes_ {\mathbf {A}} \mathbf {Q}) \tag {5}
\]

(the tensor products being given the graduations defined in no. 5) obtained by restricting the canonical homomorphisms defined in §4, nos. 1, 2 and 4; for, if \(u: \mathbf{M} \to \operatorname{Homgr}_{\mathbf{A}}(\mathbf{N}, \mathbf{P})\) is graded of degree \(\delta\) , then, for all \(x \in \mathbf{M}_{\lambda}\) , \(u(x)\) is a graded homomorphism \(\mathbf{N} \to \mathbf{P}\) of degree \(\delta + \lambda\) and hence, for \(y \in \mathbf{N}_{\mu}\) , \(u(x)(y) \in \mathbf{P}_{\delta + \lambda + \mu}\) ; if \(v: \mathbf{M} \otimes_{\mathbf{A}} \mathbf{N} \to \mathbf{P}\) corresponds canonically to \(u\) , it is then seen that \(v\) is a graded homomorphisms of degree \(\delta\) , whence our assertion concerning (3); moreover it is seen that this homomorphism is bijective. The argument is similar for (4) and (5).

If in particular \(\mathbf{P} = \mathbf{Q} = \mathbf{A}\) in (5), then there is a canonical graded homomorphism of degree 0

\[
\mathbf {M} ^ {* \mathrm{gr}} \otimes_ {\mathbf {A}} \mathbf {N} ^ {* \mathrm{gr}} \rightarrow (\mathbf {M} \otimes_ {\mathbf {A}} \mathbf {N}) ^ {* \mathrm{gr}}.\tag{6}
\]

